\documentclass[12pt, a4paper]{article}

% --- PACOTES PRINCIPAIS ---
\usepackage[utf8]{inputenc}
\usepackage[brazil]{babel}
\usepackage[T1]{fontenc}
\usepackage{amsmath, amsfonts, amssymb}
\usepackage{geometry}
\usepackage{booktabs}
\usepackage{array}
\usepackage{graphicx}
\usepackage{caption}
\usepackage{float}
\usepackage[hidelinks]{hyperref}

% --- CONFIGURAÇÃO DAS MARGENS (ABNT) ---
\geometry{top=3cm, bottom=2cm, left=3cm, right=2cm}
\linespread{1.5}

\begin{document}

% ==========================================
% CAPA DO RELATÓRIO
% ==========================================
\begin{center}
    \textbf{\Large UNIVERSIDADE DE BRASÍLIA -- UnB}\\
    \textbf{\large Faculdade do Gama -- FGA}\\
    \vspace{0.5cm}
    \textbf{\large Disciplina: Física 1 Experimental}\\
    \vspace{2.5cm}
    {\Large \textbf{RELATÓRIO}}\\
    \vspace{0.5cm}
    {\large \textbf{Experimento III - Movimento Retilíneo Uniformemente Variado (MRUV)}}\\
    \vspace{2.5cm}
\end{center}

\begin{flushright}
    \textbf{Equipa de Trabalho:}\\
    Emanuel Matsumori Quintiliano -- Matrícula: 261040090\\
    Felipe Wermelinger da Costa Leite -- Matrícula: 261014993\\
    Rodrigo Silveira Pereira Costa -- Matrícula: 251027826\\
    Yago dos Santos Rezende -- Matrícula: 241012427\\
    \vspace{0.8cm}
    \textbf{Docente:} Prof. Dr. Wytler C. Santos\\
    \textbf{Data da Prática:} [Inserir Data]
\end{flushright}

\vfill
\begin{center}
    Brasília -- DF \\
    2026
\end{center}

\newpage

% ==========================================
% 1. INTRODUÇÃO
% ==========================================
\section{Introdução}
% TODO: Inserir a introdução teórica do experimento aqui.

% ==========================================
% 2. METODOLOGIA
% ==========================================
\section{Metodologia}

\subsection{Material Utilizado}
% TODO: Listar os materiais utilizados.

\subsection{Procedimento Experimental}
% TODO: Descrever o procedimento adotado.

% ==========================================
% 3. RESULTADOS E DISCUSSÃO
% ==========================================
\section{Resultados e Discussão}

\subsection{Dados Coletados}
% TODO: Adicionar o texto explicativo sobre a tabela.

\begin{table}[!htbp]
\centering
\caption{Valores experimentais, velocidades e acelerações para o MRUV.}
\label{tab:mruv}
\vspace{0.2cm}
\resizebox{\textwidth}{!}{%
\begin{tabular}{lccccc}
\toprule
\textbf{Parâmetro} & $x_1 = 10\text{ cm}$ & $x_2 = 20\text{ cm}$ & $x_3 = 30\text{ cm}$ & $x_4 = 40\text{ cm}$ & $x_5 = 50\text{ cm}$ \\
\midrule
$t_1$ (s) & - & - & - & - & - \\
$t_2$ (s) & - & - & - & - & - \\
$t_3$ (s) & - & - & - & - & - \\
$t_4$ (s) & - & - & - & - & - \\
$t_5$ (s) & - & - & - & - & - \\
\midrule
\textbf{Tempo médio (s)} & - & - & - & - & - \\
\textbf{Erro aleatório (s)} & - & - & - & - & - \\
\textbf{Erro absoluto $u_c(t)$ (s)} & - & - & - & - & - \\
\midrule
\textbf{Velocidade Média (cm/s)} & - & - & - & - & - \\
\textbf{Velocidade Absoluta (cm/s)} & - & - & - & - & - \\
\textbf{Aceleração (cm/s$^2$)} & - & - & - & - & - \\
\bottomrule
\end{tabular}%
}
\end{table}

\subsection{Demonstração da Propagação de Incertezas}
% TODO: Inserir a demonstração dos cálculos de incerteza aqui.


% ==========================================
% 4. CONCLUSÃO
% ==========================================
\section{Conclusão}
% TODO: Redigir a conclusão final após a inserção e análise completa dos dados e gráficos.

% ==========================================
% REFERÊNCIAS
% ==========================================
\section*{Referências}
\begin{enumerate}
    \item SANTOS, Wytler C. \textbf{Roteiro para o Experimento III - Movimento Retilíneo Uniformemente Variado (MRUV)}. Faculdade UnB Gama.
\end{enumerate}

\end{document}